\chapter{Related Work}
\label{chap:relatedwork}
In today's competitive environment, shortening product life cycles place increasing
pressure on companies to manage their resources, including workforce and time, as
efficiently as possible. As \cite{Zimmermann2006} observe, project-based work has
become the dominant organisational form across most industries, making systematic and
methodologically sound project management an essential capability for any organisation
seeking to remain competitive. Project management encompasses the full lifespan of a
project, covering the ``planning, scheduling and control of project activities to achieve
performance, cost and time objectives for a given scope of work, while using resources
efficiently and effectively'' \parencite[p.~4]{Demeulemeester2002}. A project itself is
formally defined by the ISO 10006 standard as a ``unique process undertaken to achieve
an objective'' \parencite[p.~2]{ISO2017}, which ``generally consists of a set of
coordinated and controlled activities with start and finish dates, conforming to
specific requirements, including the constraints of time, cost and resources''
\parencite[p.~2]{ISO2017}. In the project planning phase, each activity is
characterised by its duration, resource demands, associated costs, and precedence
constraints that govern the order in which activities may be executed. While several
types of precedence relationships exist in practice, this thesis focuses exclusively on
the finish-to-start relationship, in which a successor activity may only begin once its
predecessor has been fully completed.
Project scheduling, as a core component of project management, involves the
construction of a base plan that specifies for each activity its precedence- and
resource-feasible start and completion times, the quantities of each resource type
required during every time period, and the resulting budget allocation
\parencite[p.~10]{Demeulemeester2002}. Formally, the goal of project scheduling is to
determine a schedule \(S = (S_1, S_2, \ldots, S_N)\) that assigns a start time \(S_i\)
to each activity \(i \in V\), such that all precedence constraints and resource capacity
constraints are satisfied while minimising the project makespan \(C_{\max}\).
A general way to represent a project is through an Activity-on-Node (AoN) network, in
which each activity \(i\) is represented as a node and directed arcs between nodes
encode the precedence constraints between activities. Throughout this thesis, all AoN
networks are assumed to be Directed Acyclic Graphs (DAGs), ensuring that no circular
dependencies exist between activities. Formally, a project is modelled as a DAG
\(G = (V, A)\), where \(V = \{0, 1, \ldots, N, N+1\}\) denotes the set of nodes
comprising the \(N\) real activities \(\{1, \ldots, N\}\) and two dummy activities
\(\{0, N+1\}\) representing the project source and sink respectively. The source
activity \(0\) has no predecessors and marks the start of the project, while the sink
activity \(N+1\) has no successors and marks its completion. The set
\(A \subseteq V \times V\) contains the directed arcs that define the precedence
relationships between activities. This thesis adopts the finish-to-start relationship
exclusively, under which arc \((i, j) \in A\) implies that activity \(j\) may not begin
before activity \(i\) has been fully completed, formally:
\[
  S_j \geq S_i + d_i \quad \forall\, (i, j) \in A
\]
where \(S_i\) denotes the start time of activity \(i\) and \(d_i\) its duration.
This chapter introduces the project scheduling problems that form the foundation of this
thesis. Section~\ref{sec:rcpsp} begins with a formal description of the most fundamental
variant, the Resource-Constrained Project Scheduling Problem (RCPSP), presenting an
integer programming formulation with a detailed explanation of each constraint, followed
by a review of the solution methods developed for this problem.
Section~\ref{sec:srcpsp} then introduces the Stochastic Resource-Constrained Project
Scheduling Problem (SRCPSP) as a natural extension that incorporates uncertainty in
activity durations. Section~\ref{sec:frcpsp} subsequently presents the Flexible
Stochastic Resource-Constrained Project Scheduling Problem (FSRCPSP), which further
extends the SRCPSP by allowing flexible resource allocation profiles, which is the core problem
addressed in this thesis. Section~\ref{sec:ml_co} reviews the emerging literature on
machine learning for combinatorial optimisation, which underpins the group-based
decomposition developed in Chapter~\ref{chap:grouprf}. The RCPSP section is organised around exact, heuristic, metaheuristic, and hybrid methods. The stochastic and flexible sections instead focus on policy classes,
scenario-based optimisation, and flexible resource-allocation models. The chapter
concludes by identifying the key limitation shared by existing approaches, namely the fixed,
non-adaptive decomposition of the search space, which motivates the hybrid approach
developed in this thesis.
% ============================================================
\section{Resource-Constrained Project Scheduling Problem}
\label{sec:rcpsp}
% ============================================================
As a basic problem in project scheduling, the deterministic RCPSP is ``probably the
most studied optimisation problem in project scheduling''
\parencite[p.~4]{SchwindtZimmermann2015}. The objective is to determine an optimal
schedule \(S^{*}\) that minimises the project makespan while adhering to precedence and
resource constraints. \cite{Blazewicz1983} showed that the RCPSP is NP-hard, using the
\(\alpha|\beta|\gamma\) notation introduced to the topic by \cite{Brucker1999}. The
problem is denoted as \(PS|prec|C_{max}\) in standard scheduling notation, where
\textit{PS} refers to project scheduling, \textit{prec} indicates precedence
relationships between activities, and \(C_{max}\) is the objective of minimising the
makespan \parencite[p.~29]{Schramme2014}.
\begin{figure}[h]
  \centering
  \includegraphics[width=0.8\textwidth]{chapters/RCPSP.png}
  \caption{RCPSP example given a single resource per activity}
  \label{RCPSP}
\end{figure}
Figure~\ref{RCPSP} illustrates the RCPSP. Each activity $i \in \{1, \dots, N\}$ has a
deterministic duration $d_i$ and a resource demand $r_i$. The figure omits the resource
index $k$, as only one resource is considered. Resources are renewable, and are therefore
replenished to their full capacity $R_{k}^{\text{max}}$ in each time period.
An early integer programming formulation of the problem is given by \cite{Pritsker1969},
who introduce the binary decision variable
\[
  x_{it} = \begin{cases}
    1, & \text{if activity } i \text{ starts at time } t \\
    0, & \text{otherwise}
  \end{cases}
\]
which determines the start time of each activity \(i\).
\begin{alignat}{2}
  & \text{Minimise:} & \quad & \sum_{t=0}^{T} t \cdot x_{N+1,t} \label{eq:rcpsp_obj} \\
  & \text{s.t.}      &       & \sum_{t=0}^{T} x_{it} = 1
    \quad \forall i \in V \label{eq:rcpsp_once} \\
  & & & \sum_{t=0}^{T} t \cdot x_{j,t} - \sum_{t=0}^{T} t \cdot x_{i,t} \geq d_j
    \quad \forall (i,j) \in A \label{eq:rcpsp_prec} \\
  & & & \sum_{i=0}^{N+1} \sum_{\tau = t - d_i + 1}^{t} r_{ik} \cdot x_{i\tau} \leq R_k
    \quad \forall t \in \{0,\ldots,T\},\, \forall k \in R \label{eq:rcpsp_res} \\
  & & & x_{it} \in \{0,1\}
    \quad \forall i \in V,\, t \in \{0,\ldots,T\} \label{eq:rcpsp_bin}
\end{alignat}
Equation~\eqref{eq:rcpsp_obj} defines the objective: schedule the sink activity as
early as possible so that the total project makespan is minimised.
Equation~\eqref{eq:rcpsp_once} enforces the non-preemptive assumption: each activity
must start exactly once and cannot be interrupted or restarted.
Equation~\eqref{eq:rcpsp_prec} encodes the finish-to-start precedence constraints,
ensuring that activity $j$ starts earliest after activity $i$ has been fully completed.
Equation~\eqref{eq:rcpsp_res} ensures that at every time period $t$, the total resource
usage does not exceed the available capacity for each resource $k \in R$.
Finally, Equation~\eqref{eq:rcpsp_bin} restricts the decision variables to binary values.
\subsection*{Exact Methods}
Since 1969, many solution methods for the RCPSP have been developed. In the early
years, researchers primarily employed exact methods. Several zero-one programming
formulations were proposed, including the seminal formulation of \cite{Pritsker1969},
and subsequent extensions by \cite{Kaplan1988} and \cite{Mingozzi1998}, each
introducing different decision variable structures. For an overview of these approaches,
the reader is referred to \cite[p.~209 ff.]{Demeulemeester2002} and the respective
papers.
Another widely used class of exact methods is Branch-and-Bound (BnB). As noted by
\cite{Demeulemeester2002}, BnB is ``the only technique which allows for the generation
of optimal solutions within an acceptable computational effort''
\parencite[p.~219]{Demeulemeester2002}. The algorithm systematically explores the
solution space by recursively decomposing the original problem into smaller subproblems
(the \textit{branching} step) while eliminating subproblems that cannot yield a
better solution than the current best (the \textit{bounding} step). This produces a
search tree in which the root node represents the full problem and each child node
represents a constrained subproblem. At each node, a lower bound on the optimal
objective value is computed; if this bound is no better than the incumbent solution, the
node is pruned and its subtree is discarded. The algorithm terminates either when all
unpruned nodes have been examined, which guarantees optimality, or when the gap between
the best found solution and the lower bound falls below a predefined threshold
\parencite[p.~219 ff.]{Demeulemeester2002}.
Building on these foundations, numerous BnB procedures have been developed, as surveyed
in \cite{Herroelen1992} and \cite{Brucker1999}. Earlier BnB procedures for the RCPSP are surveyed by \textcite{Herroelen1992}. Later, \textcite{Brucker1998} introduced a schedule-scheme representation that provides a compact way to represent sets of feasible schedules and supports systematic search tree
enumeration. Despite these advances, exact methods face a fundamental scalability barrier. As \cite{Pellerin2020} observe, ``the best exact methods can only solve instances involving at most 60 activities where
the instances are not highly resource-constrained'' (p.~397). Since real-world projects
regularly involve far more activities and tighter resource constraints, the research
community has progressively shifted its focus toward heuristic and metaheuristic
approaches.
\subsection*{Schedule Generation Schemes and Priority Rules}
Heuristic approaches for the RCPSP are primarily built upon two components:
\textit{Schedule Generation Schemes} (SGS) and \textit{priority rules}. SGS form the
foundation of most heuristic RCPSP procedures \parencite{Kolisch1999} and determine how
a feasible schedule is constructed from a given activity ordering. The two fundamental
SGS variants are the \textit{Serial SGS} (SSGS), originally introduced by
\cite{Kelley1963}, and the \textit{Parallel SGS} (PSGS), introduced by
\cite{Brooks1965}. The key distinction lies in their incrementation strategy: the SSGS
performs \textit{activity-incrementation}, scheduling one activity at a time in a fixed
priority order, while the PSGS performs \textit{time-incrementation}, advancing through
discrete time steps and scheduling all eligible activities at each step
\parencite{Kolisch1996}. For a comprehensive treatment of both schemes, the reader is
referred to \cite{Kolisch1996} and \cite{Demeulemeester2002}.
Priority rules determine the order in which activities are considered for scheduling
within an SGS. They were extensively studied by \cite{Klein2000} and classified into
five categories by \cite{lawrence1985}, as summarised in \cite{Demeulemeester2002}.
Heuristics based on priority rules are further divided into \textit{single-pass}
methods, which apply one SGS with one priority rule to generate a single schedule, and
\textit{multi-pass} methods, which generate multiple schedules through techniques such
as random sampling, forward-backward scheduling, or systematic enumeration of multiple
priority rules to explore a broader portion of the solution space
\parencite[p.~152]{Kolisch1999}. For a comprehensive review of heuristic methods for
the RCPSP, the reader is referred to \cite{Kolisch1999}, \cite{hartmann2000},
\cite{Kolisch2006a}, and \cite{issa2020}.
\subsection*{Metaheuristic and Hybrid Methods}
Over the past decades, various metaheuristic approaches have been developed, most
notably Tabu Search (TS) \parencite{Glover1989,Glover1990}, Simulated Annealing (SA)
\parencite{Kirkpatrick1983}, and Genetic Algorithms (GA) \parencite{Holland1992}. TS
and SA are both classified as local search methods: TS follows a steepest-descent
strategy, systematically selecting the best move within the neighbourhood of the current
solution, whereas SA adopts a stochastic descent mechanism inspired by the physical
annealing process. GA, in contrast, is population-based and derived from the principles
of biological evolution, maintaining a population of candidate solutions that is
iteratively evolved to generate new solutions. \cite{Bouleimen2003} propose an
SA-based approach that relies on priority lists in combination with the SSGS.
\cite{Hartmann1998} introduces three GA variants based on priority-list, priority-rule,
and random-key representations, all of which employ the SSGS and a two-point crossover
operator to generate offspring. For a comprehensive overview, the reader is referred to
\cite{hartmann2000}, \cite{Kolisch2006a}, and \cite{Demeulemeester2002}.
As individual metaheuristics have reached a performance plateau,
research has shifted toward hybrid metaheuristics. These are classified as \textit{integrative
hybrids}, which combine different search paradigms (e.g., local and population-based
search), or \textit{collaborative hybrids}, which execute multiple metaheuristics
sequentially or in parallel while exchanging information
\parencite{Pellerin2020}.
% ============================================================
\section{Stochastic Resource-Constrained Project Scheduling Problem}
\label{sec:srcpsp}
% ============================================================
The deterministic RCPSP, as introduced in Section~\ref{sec:rcpsp}, is built on the
assumption that both resource demands $r_{ik}$ and activity durations $d_i$ are known
and fixed before the project begins. While this assumption makes the problem
mathematically tractable, it does not always reflect the reality of how projects unfold
in practice. In many real-world scenarios, accurately determining all activity
parameters in advance is either very difficult or entirely impossible, especially for
projects that have never been undertaken before, where no historical data or prior
experience exists to make reliable estimates. Adding to this challenge, the increasingly
competitive business environment puts companies under constant pressure to deliver
projects faster, compressing activity durations to their absolute minimum and leaving
very little room for error. If even a single activity takes longer than planned, the
entire schedule can be disrupted, potentially causing missed
deadlines and financial penalties. To address these
shortcomings of the deterministic model, \cite{Mohring1984} proposed the SRCPSP, where
activity durations are no longer treated as fixed constants but instead modelled as
random variables $\mathbf{d}_i$, drawn from a joint probability distribution
$\mathbf{P}$.

\begin{figure}[h]
  \centering
  \includegraphics[width=0.8\textwidth]{chapters/SRCPSP.png}
  \caption{SRCPSP example given a single resource per activity. Unlike the deterministic
  RCPSP in Figure~\ref{RCPSP}, each activity duration $d_i$ is a random variable drawn
  from a probability distribution, indicated by the distribution annotations on each
  node.}
  \label{SRCPSP}
\end{figure}

Analogous to the AoN representation of the RCPSP, the SRCPSP retains the same
precedence and resource constraints but replaces deterministic activity durations with
random variables, where $d = (d_1, \ldots, d_N) \sim P$ is drawn from a joint
probability distribution \parencite[p.~12]{Stork2001}. Since the actual duration of an
activity is only revealed upon its completion, a fixed deterministic schedule is no
longer meaningful. This fundamental shift in information structure necessitates a
different solution concept: instead of computing a single schedule $S$ before the
project begins, the SRCPSP requires a \textit{scheduling policy} $\pi$, i.e.\ an online
decision process that determines which activities to start at each decision point $t$,
based on the observed project state \parencite{Demeulemeester2002}. Applying policy
$\pi$ to a duration realisation $d$ yields a schedule
$\pi(d) = (S_1, \ldots, S_N)$ with makespan $C_{\max}(\pi(d))$, and the standard
objective is to minimise the expected project duration $\mathbb{E}[C_{\max}(\pi(d))]$
\parencite[p.~292]{Herroelen2005}.

\subsection*{Scheduling Policy Classes}
The shift from deterministic schedules to scheduling policies introduces a design space
of its own. \cite{Stork2001} distinguishes between several policy classes based on the
type of information they use to make scheduling decisions. The simplest class is the
\textit{activity-based} policy, which pre-computes a fixed activity list before
execution and schedules activities in that order as resources become available; this is
equivalent to applying the SSGS with a fixed priority list at runtime. While easy to
compute, activity-based policies are non-anticipative only in a limited sense: they
cannot adapt the ordering itself in response to observed durations.

A more expressive class is the \textit{resource-based} policy, which maintains a set of
priority rules or conditional orderings that depend on the project state at each
decision point. \cite{Stork2001} showed that resource-based policies dominate
activity-based policies in terms of the expected makespan they can achieve, since they
can adapt their decisions to observed duration realisations. However, this expressiveness
comes at a computational cost: optimising over the space of resource-based policies is
substantially harder than optimising a fixed activity list
\parencite[p.~65 ff.]{Stork2001}.

A third class, the \textit{pre-selective} policy, works by selecting a subset of
time-feasible schedules before execution and choosing among them at runtime based on
observed durations. \cite{Lamas2016} demonstrate that pre-selective policies can be
approximated effectively through scenario-based optimisation, connecting the policy
design problem to the Sample Average Approximation (SAA) framework that this thesis
also employs.

\subsection*{Proactive and Reactive Approaches}
Solution methods for the SRCPSP are broadly divided into proactive and reactive
strategies \parencite{Herroelen2005}. Proactive methods construct a baseline schedule
before the project begins that explicitly accounts for potential disruptions.
\cite{Herroelen2004a} proposed a proactive approach based on inserting time buffers
between activities proportional to the variance of their duration distributions. The
rationale is that activities whose durations are highly uncertain should be given more
slack in the baseline schedule, absorbing likely delays without triggering a cascade of
rescheduling. \cite{VanDeVonder2008} extended this line of work by developing
buffer-sizing heuristics that maximise schedule stability, measured as the expected
deviation between the baseline and the realised schedule, subject to a makespan
deadline.

Reactive methods, by contrast, do not commit to a baseline schedule and instead make
scheduling decisions during project execution as uncertainty is resolved.
\cite{VanDeVonder2005} showed that purely reactive approaches can outperform proactive
ones when disruptions are frequent and severe, since they avoid the cost of protecting
a baseline that may be rendered irrelevant by early disruptions. However, the absence of
a baseline makes it difficult to coordinate resources and communicate plans to
stakeholders, which limits the practical applicability of purely reactive methods in
many organisational settings \parencite{Herroelen2005}.

A middle ground is the \textit{proactive-reactive} approach, which constructs a robust
baseline schedule proactively and then applies reactive adjustments during execution
when the baseline is violated. \cite{VanDeVonder2008} found that this combined strategy
consistently outperforms either extreme alone, particularly when the degree of
uncertainty is moderate. The present thesis follows a similar philosophy:
a baseline schedule is constructed proactively via scenario-based optimisation, and
the Relax-and-Fix procedure then refines it in a structured, iterative manner.

\subsection*{Scenario-Based and Sampling Methods}
A central computational challenge in the SRCPSP is evaluating the expected makespan
$\mathbb{E}[C_{\max}(\pi(d))]$, which requires integration over the joint duration
distribution $P$. Since this integral is generally intractable for realistic instance
sizes, most practical methods approximate it through sampling. \cite{Lamas2016}
formalised the use of Sample Average Approximation (SAA) for the SRCPSP, replacing the
expectation with an empirical average over a finite set of $|S|$ sampled scenarios:
\begin{equation}
  \mathbb{E}[C_{\max}(\pi(d))]
  \approx \frac{1}{|S|} \sum_{\pi \in S} C_{\max}(\pi(d^\pi))
  \label{eq:saa_approx}
\end{equation}
The SAA approach converts the stochastic problem into a deterministic MIP that can be
solved with standard solvers, at the cost of introducing a sampling error that decreases
as $|S|$ grows. \cite{Lamas2016} provided convergence guarantees showing that the SAA
solution converges to the true optimal policy as $|S| \to \infty$, and demonstrated
empirically that moderate sample sizes ($|S| \approx 10$--$30$) produce near-optimal
solutions for PSPLIB instances. This SAA framework is adopted directly in the FSRCPSP
formulation of this thesis (see Chapter~\ref{chap:baseline}).

\cite{Ballestin2007} took a different sampling approach, embedding Monte Carlo
simulation within a GA to evaluate candidate schedules. Each individual in the GA
population is decoded into a schedule, simulated across multiple duration realisations,
and evaluated by its average makespan. This simulation-based fitness evaluation captures
the stochastic objective more naturally than deterministic surrogate objectives but
incurs a substantial computational overhead that limits population size and generation
count. The trade-off between evaluation fidelity and computational budget is a recurring
theme in stochastic scheduling methods. The decomposition-based
approach pursued in this thesis sidesteps it by having the MIP solver optimise each subproblem exactly
under the SAA approximation rather than relying on simulation.

% ============================================================
\section{Flexible Stochastic Resource-Constrained \\ Project Scheduling Problem}
\label{sec:frcpsp}
% ============================================================
The SRCPSP models uncertainty in activity durations but retains a critical assumption
from the deterministic RCPSP: resource demands $r_{ik}$ are fixed parameters,
determined before the project begins. In practice, however, the amount of resource
allocated to an activity is often itself a decision variable. A software development
task, for instance, can be staffed with two or four engineers; more resources shorten
the duration but increase the per-period cost and reduce availability for other
activities. The Flexible Resource-Constrained Project Scheduling Problem (FRCPSP)
captures this trade-off by allowing the scheduler to choose not only \textit{when}
each activity is performed but also \textit{how many resources} are allocated to it in
each time period \parencite{Naber2014}.

The foundations of this research stream were laid by \cite{Kolisch2003}, who introduced
an initial MIP formulation for the FRCPSP with discrete resources in the context of
real-world pharmaceutical research projects. Their formulation couples the scheduling
and resource allocation decisions but relies on priority rule-based heuristics,
specifically serial and parallel SGS variants, that schedule each activity as early as
possible while allocating the maximum feasible resource amount. While effective for the
pharmaceutical instances studied, this greedy allocation strategy does not explore the
full space of feasible resource profiles and can produce suboptimal schedules when
resource sharing between concurrent activities is beneficial.

\cite{Naber2014} substantially advanced the modelling framework by specifying four
time-discrete MIP formulations for the continuous-resource FRCPSP, each differing in
how the relationship between resource allocation and activity duration is represented.
They also introduced a taxonomy of three resource categories (fully flexible, partially
flexible, and fixed) that has since become a standard classification in the literature.
Crucially, \cite{Naber2014} demonstrated that the choice of MIP formulation has a
significant impact on solver performance: formulations with tighter LP relaxations
converge faster, even when they contain more variables, a finding that directly
motivates the emphasis on LP relaxation quality in the Relax-and-Fix procedure of
this thesis.

On the metaheuristic side, \cite{Tritschler2017} proposed a hybrid approach that
combines a GA with a Variable Neighbourhood Search (VNS). The GA uses the
priority-list representation of \cite{Hartmann1998} together with the non-greedy
resource allocation strategy of \cite{Fundeling2010}, which allows the GA to explore
resource profiles beyond the greedy maximum. The VNS component then refines
GA-generated solutions by transferring resources between activities in a systematic
neighbourhood search. This two-phase design reflects a broader pattern in the FRCPSP
literature: the scheduling and resource allocation decisions are sufficiently coupled
that single-method approaches struggle to handle both effectively, and hybrid
architectures that assign different mechanisms to different aspects of the problem tend
to outperform monolithic solvers.

\cite{Ranjbar2010} adopted the formulation of \cite{Kolisch2006b} to study the RCPSP
with flexible work profiles and applied a GA as the solution method, confirming the
effectiveness of evolutionary approaches for this problem class but also highlighting
their sensitivity to representation design: the encoding of resource profiles within
the GA chromosome significantly affects convergence speed and solution quality.
\cite{Naber2017} addressed the continuous-time variant of the FRCPSP using a
branch-and-cut approach, providing both a comprehensive taxonomy of FRCPSP publications
and strong computational results on medium-sized instances. Separately,
\cite{Kellenbrink2015} extended the multi-mode RCPSP of \cite{Talbot1982} by
introducing a model-endogenous decision framework for flexible project structures,
applying a GA to the combined activity selection and scheduling problem.

What these works share, however, is a limitation: existing FRCPSP
approaches are predominantly deterministic, and existing SRCPSP approaches assume fixed
resource profiles. These two sources of complexity, stochastic workloads and flexible
resource allocation, have been studied in isolation but not in combination. The
Flexible Stochastic Resource-Constrained Project Scheduling Problem (FSRCPSP)
formulated in this thesis addresses this gap by incorporating both within a single
model, using a chance-constrained (CC) formulation approximated
via Sample Average Approximation (SAA). Rather than employing evolutionary algorithms as in the FRCPSP
literature, the solution approach embeds machine learning directly within a
Relax-and-Fix decomposition framework, using feature-based activity grouping to guide
the MIP subproblem structure, a direction not previously explored in either the FRCPSP
or SRCPSP literature.

% ============================================================
\section{Machine Learning for Combinatorial Optimisation}
\label{sec:ml_co}
% ============================================================
Over the past decade, machine learning has increasingly been used to guide and accelerate combinatorial optimisation. The underlying motivation is that classical solvers often fail to exploit problem-specific structure that could, in principle, be
learned from data. \cite{Bengio2021} provide a comprehensive survey, broadly
classifying approaches into three categories: \textit{end-to-end learning} methods that
replace the solver entirely with a neural network, \textit{configuration} methods that
learn solver parameters or decomposition structures from training instances, and
\textit{hybrid} methods that combine learned components with exact or heuristic search
while preserving the solver's optimality guarantees. The group-based approach proposed
in this thesis falls within the hybrid category: an unsupervised learning pipeline
determines the decomposition structure (i.e.\ the activity grouping), while the MIP solver
retains full responsibility for solving each resulting subproblem.

\subsection*{Learning to Decompose and Configure MIP Solvers}
Much of the ML-for-CO literature concerns learning \textit{how} to
decompose a problem rather than learning to solve it outright. \cite{Khalil2016}
introduced a framework in which a machine learning model learns to select branching
variables in branch-and-bound, replacing the hand-crafted heuristics that commercial
solvers use by default. Their approach trains a classifier on features extracted from
the LP relaxation at each node (including variable fractionality, objective
coefficient, and constraint participation) and learns a branching policy that
generalises across instances from the same problem family. Branching decisions, while individually simple, collectively determine the shape of the
search tree, and a learned policy can prune the tree far more aggressively than generic
rules. The activity grouping pipeline in this thesis operates on a similar principle:
it determines which binary variables are active in each subproblem, thereby
shaping the LP relaxation and branching behaviour within that subproblem.

\cite{Lodi2020} survey the broader landscape of ML-assisted MIP solving, identifying
three levels at which learning can intervene: algorithm selection (choosing which solver
or configuration to use), algorithm configuration (tuning solver parameters for a
problem class), and algorithm design (replacing solver components with learned
alternatives). The decomposition strategy pursued in this thesis operates at the
algorithm design level: the grouping pipeline replaces a fixed activity-by-activity
decomposition with a learned, instance-adaptive one.

\subsection*{Neural Constructive and Improvement Methods}
At the other end of the spectrum, several approaches aim to replace the solver entirely
with a learned construction or improvement procedure. \cite{Vinyals2015} introduced the
Pointer Network, a sequence-to-sequence model with an attention mechanism that learns
to output permutations, which is a natural fit for routing problems where the solution is an
ordering of nodes. \cite{Kool2019} extended this approach with a more efficient
attention-based architecture and reinforcement learning training, achieving
near-optimal results on the Travelling Salesman Problem for instances up to 100 nodes.

For improvement rather than construction, \cite{Hottung2020} proposed a neural large
neighbourhood search (LNS) in which a neural network learns a destroy-and-repair
operator: the network selects which portion of the current solution to destroy and
reconstructs it, replacing the hand-designed operators traditionally used in LNS. Their
results on vehicle routing problems demonstrate that learned operators can outperform
portfolios of classical destroy-repair heuristics, particularly when the operator is
trained on the same problem family as the target instances.

While these end-to-end methods achieve impressive results on well-studied benchmark
problems, they share a common limitation that makes them unsuitable for the FSRCPSP:
they require large training datasets of solved instances from the same problem family.
For a novel problem formulation such as the FSRCPSP, no such training set exists, and
generating one would require solving many FSRCPSP instances to optimality, which is precisely
the task the method is supposed to accelerate. The unsupervised approach taken in this
thesis avoids this chicken-and-egg problem by learning structure from a single instance
at solve time rather than from a corpus of pre-solved instances.

\subsection*{Feature Engineering for Scheduling}
Regardless of whether learning is used for construction, improvement, or decomposition,
the choice of features plays a critical role. The features must provide a minimum
sufficient representation of each decision unit, capturing the aspects of its scheduling
behaviour that are relevant for the learning task without introducing redundancy that
would inflate the feature space and impair learning quality through the curse of
dimensionality \parencite{Bellman1954, Guyon2003}.

For project scheduling specifically, feature design has been studied most extensively
in the context of priority rule development. \cite{Kolisch1999} demonstrated that
structural network features (activity slack, resource demand, in-degree, out-degree)
are among the most informative features for scheduling decisions, consistently
outperforming random or index-based orderings across diverse instance sets. More
recently, \cite{Hottung2020} and \cite{Khalil2016} have shown that features derived
from the LP relaxation (fractionality, reduced cost, constraint tightness) can further
improve learning-based scheduling decisions, though these features are specific to
methods that interact with the LP solver during search.

The feature set used in this thesis (Chapter~\ref{chap:grouprf}) draws on both
traditions: network position features (precedence level, in-degree, out-degree) capture
the structural role of each activity in the project DAG, scheduling urgency features
(slack, pressure index) capture how constrained each activity's time window is, and
stochastic resource features (mean workload, coefficient of variation) capture the
resource profile across scenarios. Together, these three groups answer distinct questions about each activity: where it sits in the network, how urgently it needs to
be scheduled, and what resource pressure it creates, while keeping
redundancy between groups low.

\subsection*{Unsupervised Learning for Problem Decomposition}
The grouping mechanism in this thesis uses Ward hierarchical clustering
\parencite{Ward1963} to partition activities into groups based on feature similarity.
Ward's method provides a principled agglomerative criterion that minimises the increase
in total within-cluster variance at each merge step, producing compact and
well-separated groups from heterogeneous feature vectors. The number of clusters is
selected via the silhouette coefficient \parencite{Rousseeuw1987}, which measures the
ratio of within-cluster cohesion to between-cluster separation and enables data-driven
cut selection without manual tuning.

The use of unsupervised clustering to guide MIP decomposition is relatively unexplored
in the operations research literature. Existing decomposition strategies, such as the
product-, resource-, and process-oriented decompositions studied by \cite{Helber2010}
for lot-sizing, or the unit-oriented decomposition of \cite{bigler2024multisite} for
multi-site project scheduling, rely on domain knowledge to define the decomposition
structure. The approach proposed in this thesis replaces this manual design with a
data-driven alternative that adapts the decomposition to each instance's specific
structure. To the authors' knowledge,
no prior work has applied clustering-guided decomposition to project scheduling problems
under stochastic, flexible resources, making this a methodological contribution at the intersection of
unsupervised learning and MIP-based scheduling.
